\documentclass[11pt,a4paper]{article}
\usepackage[italian]{babel}
\usepackage[margin=2.5cm]{geometry}
\usepackage{parskip}
\usepackage{amsmath, amssymb, amsthm}
\usepackage{booktabs}
\usepackage{array}
\usepackage{xcolor}
\usepackage{needspace}
\usepackage{enumitem}
\usepackage{listings}
\usepackage{tikz}
\usepackage[most]{tcolorbox}
\usepackage[hidelinks]{hyperref}
\usetikzlibrary{decorations.pathreplacing, shapes.geometric, arrows.meta, positioning}

% Niente vedove/orfane: i paragrafi non lasciano righe isolate a fine/inizio pagina
\widowpenalty=10000
\clubpenalty=10000

% Elenchi più compatti
\setlist{itemsep=2pt, parsep=0pt, topsep=4pt, partopsep=0pt}

% --- Riquadri con bordo nero, senza numero (si spezzano tra due pagine) ---
% Uso: \begin{definizione} ... \end{definizione}
%      \begin{teorema}[Nome] ... \end{teorema}  (il nome è facoltativo)
\tcbset{nota/.style={enhanced jigsaw, breakable, colback=white,
  colframe=black, boxrule=0.5pt, sharp corners}}
\NewTColorBox{definizione}{}{nota,
  before upper={\textbf{Definizione.}\ }}
\NewTColorBox{teorema}{o}{nota,
  before upper={\textbf{Teorema\IfValueT{#1}{ (#1)}.}\ }}
\NewTColorBox{proposizione}{o}{nota,
  before upper={\textbf{Proposizione\IfValueT{#1}{ (#1)}.}\ }}
\NewTColorBox{proprieta}{o}{nota,
  before upper={\textbf{Proprietà\IfValueT{#1}{ (#1)}.}\ }}
\NewTColorBox{assioma}{o}{nota, boxrule=1pt,
  before upper={\textbf{Assioma\IfValueT{#1}{ (#1)}.}\ }}

% --- Ambienti semplici (senza riquadro) ---
% Il titolo sta in cima al blocco, anche se il contenuto inizia con un riquadro o
% con codice; e non resta mai da solo in fondo a una pagina.
% Uso: \begin{esempio}[Titolo facoltativo] ... \end{esempio}
\newcommand{\newrunin}[2]{%
  \NewDocumentEnvironment{#1}{o}{%
    \par\Needspace{4\baselineskip}\addvspace{\medskipamount}\noindent
    \textbf{#2}\IfValueT{##1}{ (##1)}\textbf{.}\ \ignorespaces}%
  {\par\addvspace{\medskipamount}}}
\newrunin{esempio}{Esempio}
\newrunin{esercizio}{Esercizio}
\newrunin{osservazione}{Osservazione}

% V e F nelle tavole di verità
\newcommand{\V}{\textsf{V}}
\newcommand{\F}{\textsf{F}}

% --- Codice C ---
% Uso: \begin{codice} ... \end{codice}      codice C numerato
%      \begin{comando} ... \end{comando}    comandi da terminale
%      \begin{risultato} ... \end{risultato} output di un programma
%      \lstinline|printf("ciao");|           codice dentro al testo
\lstdefinestyle{c}{
  language=C,
  basicstyle=\ttfamily\small,
  keywordstyle=\color{blue}\bfseries,
  commentstyle=\color{gray}\itshape,
  stringstyle=\color{red!70!black},
  numbers=left,
  numberstyle=\tiny\color{gray},
  numbersep=8pt,
  columns=flexible,
  keepspaces=true,
  breaklines=true,
  showstringspaces=false,
  tabsize=4,
  morekeywords={include, define},
  % lettere accentate nei commenti
  literate={à}{{\`a}}1 {è}{{\`e}}1 {é}{{\'e}}1 {ì}{{\`i}}1
           {ò}{{\`o}}1 {ù}{{\`u}}1 {È}{{\`E}}1 {À}{{\`A}}1
}
\lstset{style=c}
\newtcblisting{codice}{listing only, nota, left=6mm,
  listing options={style=c}}
\newtcblisting{comando}{listing only, nota,
  listing options={basicstyle=\ttfamily\small, numbers=none, language={},
    columns=flexible, keepspaces=true, breaklines=true}}
\newtcblisting{risultato}{listing only, enhanced jigsaw, breakable,
  colback=black!5, colframe=black, boxrule=0.5pt, sharp corners,
  listing options={basicstyle=\ttfamily\small, numbers=none, language={},
    columns=flexible, keepspaces=true, breaklines=true}}

% --- Diagrammi di flusso (simboli standard) ---
\tikzset{
  terminale/.style = {ellipse, draw, minimum width=2.5cm, minimum height=0.9cm},
  io/.style        = {trapezium, trapezium left angle=70, trapezium right angle=110,
                      draw, minimum width=2.5cm, minimum height=0.9cm},
  processo/.style  = {rectangle, draw, minimum width=2.5cm, minimum height=0.9cm},
  decisione/.style = {diamond, draw, aspect=2, minimum width=2.5cm},
  freccia/.style   = {thick, -{Stealth}}
}

\title{Logica}
\author{Marco Cerbella}
\date{Lezione 1 -- 24 settembre 2026}

\begin{document}
\maketitle

\section{Il calcolo delle proposizioni}

\begin{definizione}
Si chiama \emph{proposizione} un enunciato di cui si può stabilire con
certezza se è vero o falso.
\end{definizione}

\begin{esempio}
``$7$ è un numero primo'' è una proposizione (vera);
``$4 > 9$'' è una proposizione (falsa);
``che bella giornata!'' e ``$x > 3$'' \textbf{non} sono proposizioni:
la prima non è né vera né falsa, la seconda dipende da $x$.
\end{esempio}

Le proposizioni si indicano con lettere maiuscole $P, Q, R, \dots$ e il loro
\emph{valore di verità} è \V\ (vero) oppure \F\ (falso).

I concetti che non si definiscono, perché si ritengono noti, vengono detti
\emph{primitivi} (ad esempio ``insieme'', ``elemento'', ``vero'', ``falso'').

La logica classica si basa su due principi:
\begin{enumerate}
    \item \textbf{principio di non contraddizione}: una proposizione non può
          essere contemporaneamente vera e falsa;
    \item \textbf{principio del terzo escluso}: una proposizione è vera
          oppure falsa, non esiste una terza possibilità.
\end{enumerate}

\section{Proposizioni composte e connettivi logici}

Le proposizioni semplici si combinano tramite i \emph{connettivi logici}
per formare proposizioni composte. Il valore di verità di una proposizione
composta dipende solo dai valori delle proposizioni che la compongono, e si
riassume in una \emph{tavola di verità}.

\subsection{Negazione $\neg$}

\begin{definizione}
La \emph{negazione} di $P$, indicata con $\neg P$ (``non $P$''), è vera
quando $P$ è falsa e falsa quando $P$ è vera.
\end{definizione}

\begin{center}
\begin{tabular}{c|c}
$P$ & $\neg P$ \\ \hline
\V & \F \\
\F & \V
\end{tabular}
\end{center}

\begin{esempio}
$P$: ``$5$ è pari'' (\F). \quad $\neg P$: ``$5$ non è pari'' (\V).
\end{esempio}

Vale $\neg(\neg P) = P$: negare due volte riporta alla proposizione di
partenza.

\subsection{Congiunzione $\land$ e disgiunzione $\lor$}

\begin{definizione}
La \emph{congiunzione} $P \land Q$ (``$P$ e $Q$'') è vera solo se $P$ e $Q$
sono entrambe vere.

La \emph{disgiunzione inclusiva} $P \lor Q$ (``$P$ o $Q$'') è vera se
almeno una tra $P$ e $Q$ è vera; è falsa solo se sono entrambe false.
\end{definizione}

\begin{center}
\begin{tabular}{cc|c|c}
$P$ & $Q$ & $P \land Q$ & $P \lor Q$ \\ \hline
\V & \V & \V & \V \\
\V & \F & \F & \V \\
\F & \V & \F & \V \\
\F & \F & \F & \F
\end{tabular}
\end{center}

\begin{osservazione}
Il ``$\lor$'' è \emph{inclusivo} (come il latino \emph{vel}): è vero anche
quando $P$ e $Q$ sono vere tutte e due. Nel linguaggio comune ``o'' è spesso
esclusivo (``o mangi o esci''), ma in matematica si intende sempre quello
inclusivo, salvo indicazione contraria.
\end{osservazione}

\textbf{Collegamento con gli insiemi.}
$\land$ corrisponde all'intersezione, $\lor$ all'unione e $\neg$ al
complementare:
\[
x \in A \cap B \iff (x \in A) \land (x \in B), \qquad
x \in A \cup B \iff (x \in A) \lor (x \in B).
\]

\begin{teorema}[Leggi di De Morgan]
\[
\neg(P \land Q) \iff \neg P \lor \neg Q, \qquad
\neg(P \lor Q) \iff \neg P \land \neg Q.
\]
\end{teorema}

Sono le stesse leggi viste per gli insiemi: la negazione ``entra'' nella
parentesi e $\land$ e $\lor$ si scambiano.

\begin{esempio}
La negazione di ``piove \textbf{e} fa freddo'' è ``non piove \textbf{o} non
fa freddo'' (basta che una delle due sia falsa).
\end{esempio}

\subsection{Implicazione $\Rightarrow$}

\begin{definizione}
L'\emph{implicazione} $P \Rightarrow Q$ (``se $P$ allora $Q$'', ``$P$ implica
$Q$'') è falsa \textbf{solo} quando $P$ è vera e $Q$ è falsa. In tutti gli
altri casi è vera. $P$ si chiama \emph{ipotesi} (o antecedente), $Q$
\emph{tesi} (o conseguente).
\end{definizione}

\begin{center}
\begin{tabular}{cc|c}
$P$ & $Q$ & $P \Rightarrow Q$ \\ \hline
\V & \V & \V \\
\V & \F & \F \\
\F & \V & \V \\
\F & \F & \V
\end{tabular}
\end{center}

\textbf{Perché è vera quando $P$ è falsa?} Pensa a una promessa: ``se
prendi 30, ti regalo la bici''. L'unico modo in cui la promessa viene
\emph{tradita} è prendere 30 e non ricevere la bici (\V$\Rightarrow$\F).
Se non prendi 30, la promessa non è stata violata, qualunque cosa succeda
alla bici: per questo l'implicazione è vera.

\begin{esempio}
``Se $n$ è multiplo di $4$, allora $n$ è pari'' è vera per ogni $n$.
Per $n = 6$ l'ipotesi è falsa e la tesi vera; per $n = 7$ entrambe false:
in nessun caso l'implicazione viene violata.
\end{esempio}

\textbf{Condizione necessaria e sufficiente.} Se $P \Rightarrow Q$ è vera:
\begin{enumerate}
    \item $P$ è \emph{condizione sufficiente} per $Q$: basta che valga $P$
          perché valga $Q$;
    \item $Q$ è \emph{condizione necessaria} per $P$: se non vale $Q$, non
          può valere nemmeno $P$.
\end{enumerate}
Ad esempio ``essere multiplo di $4$'' è sufficiente per ``essere pari'', ed
``essere pari'' è necessario per ``essere multiplo di $4$'' (ma non
sufficiente: $6$ è pari ma non multiplo di $4$).

\textbf{Implicazioni collegate.} Data $P \Rightarrow Q$:
\begin{center}
\begin{tabular}{l|l|l}
nome & forma & equivalente a $P \Rightarrow Q$? \\ \hline
inversa (o reciproca) & $Q \Rightarrow P$ & no \\
contraria & $\neg P \Rightarrow \neg Q$ & no \\
\textbf{contronominale} & $\neg Q \Rightarrow \neg P$ & \textbf{sì}
\end{tabular}
\end{center}

\begin{teorema}[Contronominale]
$(P \Rightarrow Q) \iff (\neg Q \Rightarrow \neg P)$.
\end{teorema}

Su questo si basa la \emph{dimostrazione per assurdo}: invece di dimostrare
$P \Rightarrow Q$, si suppone che $Q$ sia falsa e si arriva a negare $P$
(o a una contraddizione). Anche la dimostrazione di $\sqrt{2} \notin
\mathbb{Q}$ funziona così.

Vale inoltre $(P \Rightarrow Q) \iff (\neg P \lor Q)$, da cui si ricava la
negazione dell'implicazione:
\[
\neg(P \Rightarrow Q) \iff P \land \neg Q.
\]
Cioè: per smentire ``se $P$ allora $Q$'' bisogna trovare un caso in cui
$P$ è vera e $Q$ è falsa.

\subsection{Doppia implicazione $\Leftrightarrow$}

\begin{definizione}
La \emph{doppia implicazione} (o \emph{equivalenza}) $P \Leftrightarrow Q$
(``$P$ se e solo se $Q$'') è vera quando $P$ e $Q$ hanno lo stesso valore
di verità, cioè sono entrambe vere o entrambe false.
\end{definizione}

\begin{center}
\begin{tabular}{cc|c}
$P$ & $Q$ & $P \Leftrightarrow Q$ \\ \hline
\V & \V & \V \\
\V & \F & \F \\
\F & \V & \F \\
\F & \F & \V
\end{tabular}
\end{center}

Equivale a $(P \Rightarrow Q) \land (Q \Rightarrow P)$: $P$ è condizione
\emph{necessaria e sufficiente} per $Q$. Per questo, per dimostrare un
``se e solo se'' si dimostrano le due implicazioni separatamente.

\begin{esempio}
``$n$ è pari $\Leftrightarrow$ $n^2$ è pari'' è vera per ogni
$n \in \mathbb{Z}$.
\end{esempio}

\subsection{Tautologie e contraddizioni}

\begin{definizione}
Una \emph{tautologia} è una proposizione composta sempre vera, qualunque
siano i valori delle proposizioni che la compongono.

Una \emph{contraddizione} è una proposizione composta sempre falsa.
\end{definizione}

\begin{esempio}
$P \lor \neg P$ è una tautologia (terzo escluso);
$P \land \neg P$ è una contraddizione (non contraddizione).
\begin{center}
\begin{tabular}{c|c|c|c}
$P$ & $\neg P$ & $P \lor \neg P$ & $P \land \neg P$ \\ \hline
\V & \F & \V & \F \\
\F & \V & \V & \F
\end{tabular}
\end{center}
\end{esempio}

Tutte le leggi viste (De Morgan, contronominale, \dots) sono tautologie: si
verificano costruendo la tavola di verità e controllando che l'ultima
colonna contenga solo \V.


\section{Esempi dalla lezione}

\begin{esempio}[Proposizioni e non proposizioni]
Sono proposizioni: ``Anna e Mario sono figli di Luigi'', ``Francesco Guccini è morto quest'anno''. Non sono proposizioni: ``che tempo fa?'' (interrogativa), ``è molto tardi'' (giudizio personale).
\end{esempio}

\begin{esempio}[Negazione]
``Roma è una città'' (\V) $\to$ ``Roma non è una città'' (\F).
``Il Po è un lago'' (\F) $\to$ ``Il Po non è un lago'' (\V).
\end{esempio}

\begin{esempio}[Congiunzione]
\begin{itemize}
    \item ``Roma è la capitale d'Italia e Roma è la sede del governo'': \V;
    \item ``$4$ è positivo e $4$ è dispari'': \F;
    \item ``$20$ è maggiore di $30$ e $20$ è multiplo di $3$'': \F.
\end{itemize}
\end{esempio}

\begin{esempio}[Disgiunzione]
\begin{itemize}
    \item ``Firenze è una città o l'Arno è un lago'': \V;
    \item ``$10$ è multiplo di $3$ o $-10$ è un numero negativo'': \V\ (basta che una delle due sia vera);
    \item ``Roma è la capitale della Francia o l'Italia è una città'': \F.
\end{itemize}
\end{esempio}

\begin{esempio}[Implicazione ed equivalenza]
``Se oggi è sabato allora domani è domenica''; ``se $a$ è divisibile per $6$ allora $a$ è divisibile per $3$''.

``Se Anna è milanese allora Anna è italiana'': vale $P \Rightarrow Q$ ma non $Q \Rightarrow P$, quindi $P$ e $Q$ \textbf{non} sono equivalenti. Invece ``un numero è divisibile per $2$ se e solo se è pari'': valgono $P \Rightarrow Q$ e $Q \Rightarrow P$, quindi $P \Leftrightarrow Q$.
\end{esempio}

\section{Predicati}

Il calcolo delle proposizioni vale solo per enunciati che hanno un valore
di verità. Un enunciato come
\[
P(x):\ \text{``$x$ è un numero maggiore di $100$''}
\]
non è una proposizione: non si può dire se è vero o falso finché non si
sceglie $x$.

\begin{definizione}
Un \emph{predicato} (o \emph{enunciato aperto}) è un enunciato $P(x)$ che
contiene una variabile $x$, appartenente a un insieme $U$ detto
\emph{dominio}, e che diventa una proposizione quando a $x$ si sostituisce
un elemento di $U$.
\end{definizione}

\begin{esempio}
Con $U = \mathbb{N}$ e $P(x)$: ``$x > 100$'': $P(150)$ è vera,
$P(3)$ è falsa.
\end{esempio}

\begin{definizione}
L'\emph{insieme di verità} di $P(x)$ è l'insieme degli elementi di $U$ che
rendono vero il predicato:
\[
\{x \in U \mid P(x) \text{ è vera}\}.
\]
\end{definizione}

È proprio la definizione di un insieme \emph{per caratteristica}. Nell'esempio
l'insieme di verità è $\{x \in \mathbb{N} \mid x > 100\} = \{101, 102, \dots\}$.

Un predicato può avere più variabili: $P(x, y)$: ``$x < y$''.

\subsection{Quantificatori}

Un predicato diventa una proposizione anche \emph{quantificando} la
variabile, cioè dicendo per quanti elementi di $U$ il predicato è vero.

\begin{definizione}
\begin{enumerate}
    \item \textbf{Quantificatore universale} $\forall$ (``per ogni''):
          $\forall x \in U,\ P(x)$ è vera se $P(x)$ è vera per
          \textbf{tutti} gli elementi di $U$.
    \item \textbf{Quantificatore esistenziale} $\exists$ (``esiste''):
          $\exists\, x \in U : P(x)$ è vera se $P(x)$ è vera per
          \textbf{almeno un} elemento di $U$.
\end{enumerate}
\end{definizione}

Si usa anche $\exists!$ (``esiste ed è unico''): $P(x)$ è vera per
\textbf{uno e un solo} elemento di $U$.

\begin{esempio}\mbox{}
\begin{enumerate}
    \item $\forall x \in \mathbb{R},\ x^2 \geq 0$ \quad è vera;
    \item $\forall x \in \mathbb{N},\ x > 100$ \quad è falsa (basta
          $x = 3$);
    \item $\exists\, x \in \mathbb{N} : x > 100$ \quad è vera (ad esempio
          $x = 150$);
    \item $\exists\, x \in \mathbb{Q} : x^2 = 2$ \quad è falsa, mentre
          $\exists\, x \in \mathbb{R} : x^2 = 2$ è vera.
\end{enumerate}
\end{esempio}

L'ultimo esempio mostra che il valore di verità dipende anche dal dominio.

\textbf{Come si verifica o si smentisce.}
\begin{center}
\begin{tabular}{l|p{5.2cm}|p{5.2cm}}
 & per dimostrare che è vera & per dimostrare che è falsa \\ \hline
$\forall x,\ P(x)$ & serve una dimostrazione valida per ogni $x$ &
basta \textbf{un} controesempio \\ \hline
$\exists\, x : P(x)$ & basta \textbf{un} esempio &
serve mostrare che $P(x)$ è falsa per ogni $x$
\end{tabular}
\end{center}

\subsection{Negazione dei quantificatori}

\begin{teorema}
\[
\neg\big(\forall x,\ P(x)\big) \iff \exists\, x : \neg P(x), \qquad
\neg\big(\exists\, x : P(x)\big) \iff \forall x,\ \neg P(x).
\]
\end{teorema}

Regola pratica: la negazione ``attraversa'' il quantificatore, che si
scambia ($\forall \leftrightarrow \exists$), e poi nega il predicato.

\begin{esempio}\mbox{}
\begin{enumerate}
    \item ``Tutti gli studenti hanno superato l'esame'' si nega con
          ``\textbf{esiste} almeno uno studente che \textbf{non} ha superato
          l'esame''. Attenzione: \textbf{non} con ``nessuno studente ha
          superato l'esame''.
    \item $\neg(\forall x \in \mathbb{R},\ x > 0) \iff
          \exists\, x \in \mathbb{R} : x \leq 0$.
\end{enumerate}
\end{esempio}

\subsection{Più quantificatori: l'ordine conta}

Con più variabili si possono usare più quantificatori. Quantificatori
\emph{uguali} si possono scambiare, quantificatori \emph{diversi} no.

\begin{esempio}
Con $U = \mathbb{N}$:
\begin{enumerate}
    \item $\forall x\ \exists\, y : y > x$ \quad (``per ogni numero ne
          esiste uno più grande'') è \textbf{vera}: dato $x$, basta
          $y = x + 1$;
    \item $\exists\, y\ \forall x : y > x$ \quad (``esiste un numero più
          grande di tutti'') è \textbf{falsa}.
\end{enumerate}
Nel primo caso $y$ può dipendere da $x$, nel secondo deve essere lo stesso
per tutti gli $x$.
\end{esempio}

Per negare più quantificatori si applica la regola uno alla volta:
\[
\neg\big(\forall x\ \exists\, y : P(x,y)\big) \iff
\exists\, x\ \forall y : \neg P(x,y).
\]
Questo tipo di negazione servirà spesso in analisi (ad esempio con le
definizioni di limite).

\end{document}
